\rhead{DS-EL1: High-Dimensional Data Analysis}
\phantomsection 
\section*{High-Dimensional Data Analysis (DS-EL1)}
\label{elec:DS_1}

\noindent \small{\hyperref[main:scheme]{$\leftarrow$ Back to Scheme}} \vspace{0.5cm}

\noindent \textbf{Credits:} 4 (3-1-0)

\subsection*{Theory}
\noindent\textbf{Unit 1} \\
\textit{High-Dimensional Phenomena and Curse of Dimensionality:} Distance concentration and emptiness phenomenon, Concentration of measure inequality, Nearest neighbor distances in high dimensions, Sparsity in high-dimensional data, Double descent phenomenon, Blessing of dimensionality vs. curse, Dimension reduction necessity, Johnson-Lindenstrauss lemma for random projections.

\vspace{0.3cm}

\noindent\textbf{Unit 2} \\
\textit{Principal Component Analysis and Classical Methods:} PCA mathematical formulation (covariance matrix eigendecomposition, SVD), PCA variants (Kernel PCA, Sparse PCA, Robust PCA), Incremental/online PCA, Factor analysis and independent component analysis (ICA), Multidimensional scaling (MDS - classical, non-metric), Isomap and geodesic distances.

\vspace{0.3cm}

\noindent\textbf{Unit 3} \\
\textit{Nonlinear Dimensionality Reduction:} t-SNE algorithm (student-t divergence, perplexity tuning), UMAP (uniform manifold approximation and projection), LargeVis and PHATE, Autoencoder architectures (vanilla, variational, denoising), Deep belief networks for representation learning, Self-supervised contrastive learning (SimCLR, MoCo), Manifold learning assumptions and topology preservation.

\vspace{0.3cm}

\noindent\textbf{Unit 4} \\
\textit{High-Dimensional Statistics and Regularization:} Multiple testing problem and FDR control (Benjamini-Hochberg procedure), Sparsity-inducing regularization (Lasso, Elastic Net, Group Lasso), Stability selection and knockoffs framework, High-dimensional covariance estimation (graphical models, covariance shrinkage), Robust high-dimensional regression, Sure independence screening (SIS).

\vspace{0.3cm}

\noindent\textbf{Unit 5} \\
\textit{Scalable Algorithms and Embeddings:} Random projection methods (CountSketch, Johnson-Lindenstrauss transforms), Locality-Sensitive Hashing (LSH) families (random hyperplanes, p-stable distributions), Approximate nearest neighbors (HNSW, FAISS), Word embeddings scaling (GloVe, fastText), Graph embeddings (Node2Vec, DeepWalk), Tensor decomposition for multi-modal data.

\newpage

\subsection*{Practical}
\noindent\textbf{Lab 1: PCA and Classical DR Comparison} \\
Focuses on eigendecomposition vs. SVD, explained variance analysis, reconstruction error.

\vspace{0.2cm}

\noindent\textbf{Lab 2: Nonlinear DR Visualization} \\
Focuses on t-SNE/UMAP parameter tuning, perplexity analysis, cluster separation metrics.

\vspace{0.2cm}

\noindent\textbf{Lab 3: Autoencoder Implementation} \\
Focuses on latent space visualization, reconstruction quality, anomaly detection use cases.

\vspace{0.2cm}

\noindent\textbf{Lab 4: High-Dimensional Regression} \\
Focuses on Lasso path algorithms, cross-validation, stability selection analysis.

\vspace{0.2cm}

\noindent\textbf{Lab 5: LSH and ANN Search} \\
Focuses on FAISS/HNSW indexing, recall@K evaluation, query time benchmarking.

\vspace{0.2cm}

\noindent\textbf{Lab 6: Random Projections Scalability} \\
Focuses on JL lemma verification, dimensionality reduction performance scaling.

\vspace{0.2cm}

\noindent\textbf{Lab 7: Graph Embeddings} \\
Focuses on Node2Vec parameter tuning, link prediction evaluation.

\vspace{0.2cm}

\noindent\textbf{Lab 8: Multiple Testing Control} \\
Focuses on FDR procedures, power analysis, volcano plots for genomics data.

\vspace{0.2cm}

\noindent\textbf{Lab 9: Sparse Covariance Estimation} \\
Focuses on graphical lasso, precision matrix visualization, conditional independence.

\vspace{0.2cm}

\noindent\textbf{Lab 10: High-Dimensional Pipeline} \\
Focuses on complete analysis workflow from raw HD data to interpretable low-D representations.

\vspace{0.2cm}

\newpage
\rhead{Curriculum Schema}
